% ========================================================================================================
% Introduction 

\minitoc 

En informatique et en mathématiques, l'optimisation est la recherche d'un optimum pour un problème posé. 
On s'intéresse ici à optimisation dite \emph{continue}. On peut alors formuler un tel problème sous la forme suivante : 

\begin{center}    
    \begin{minipage}{0.45\textwidth}
        \begin{mini*}|l|
            {}{f(x)}{}{}
            % \addConstraint{g_{j}(x)}{\le 0, j = 1, \hdots, m}
            \addConstraint{        x }{\in \mathcal{C}}
        \end{mini*}
    \end{minipage}
    \hfill 
    \begin{minipage}{0.45\textwidth}
        \begin{mini*}|s|
            {}{f(x)}{}{}
            % \addConstraint{g_{j}(x)}{\le 0, j = 1, \hdots, m}
            \addConstraint{        x }{\in \mathcal{C}}
        \end{mini*}
    \end{minipage}
\end{center}

On appelle $x$ la \emph{variable de décision} ou \emph{variable de contrôle}. La fonction $f : \R^n \longrightarrow R$ 
est appelée \emph{fonction objectif} ou le \emph{critère}. 
L'ensemble $C \subset \R$ est appelée \emph{ensemble des contraintes}. 

Résoudre un tel problème signifie trouver $ \overline{x} \in C \subset \R^n$ tel que :
    \[ \forall x \in C, \quad f(x) \geqslant \overline{x} \] 

\begin{remark}
    On peut remarque que maximiser $f$ revient à minimiser $-f$. On peut donc toujours se ramener 
    à un problème de minisation. De plus, lorsque $C = \R^n$ on dit que le problème est \emph{non contraint}. 
\end{remark}


% ========================================================================================================
% CNS d'optimalité dans contraintes

\section{CNS d'optimalité dans contraintes}

Dans cette section, nous nous intéressons aux problèmes de la forme suivante : 

\begin{mini*}|s|
    {}{f(x)}{}{}
    % \addConstraint{g_{j}(x)}{\le 0, j = 1, \hdots, m}
    \addConstraint{        x }{\in \R^n}
\end{mini*}
où $f : \R^n \longrightarrow \R$. Nous commencerons par aborder l'optimalité locale puis nous passerons, grâce 
à la convexité, à l'optimalité globale. 

\subsection{Conditions Nécessaires}

Évoquons tout d'abord quelques conditoons nécessaires d'optimalité globale. 

\begin{theorem}[CN1 - Principe de Fermat]
    Soit $f : \R^n \longrightarrow \R$ et $ \overline{x} \in \R^n$ un minimum local de $f$. 
    Si $f$ \emph{est différentiable} alors : 
        \[ \nabla f( \overline{x} ) = 0 \]  
\end{theorem}

\begin{definition}[Matrice définie positive]
    Soit $A \in \mathcal{M}_{n,m}(\R)$ une matrice, on dit que $A$ est \emph{définie positive} si 
        \[ \forall x \in \R^n, x \not = 0, \quad ^t x A x > 0  \]
        \[ \iff Sp(A) \subset R_+^* \] 
    On note alors $ A \succ 0$ 
\end{definition}

\begin{definition}[Matrice semi-définie positive]
    Soit $A \in \mathcal{M}_{n,m}(\R)$ une matrice, on dit que $A$ est \emph{semi-définie positive} si 
        \[ \forall x \in \R^n, x \not = 0, \quad ^t x A x \geqslant 0  \]
        \[ \iff Sp(A) \subset R_+ \] 
    On note alors $ A \succeq 0$ 
\end{definition}

\begin{theorem}[CN2]
    Soit $f : \R^n \longrightarrow \R$ une fonction quelconque et soit $ \overline{x} \in \R^n$ un minimum 
    local de $f$. Si $f$ est \emph{deux fois différentiable en} $ \overline{x}$ alors : 
        \[ \nabla^2 f( \overline{x}) \succeq 0 \] 
    (i.e ssi la hessienne de $f$ est semi-définie positive). 
    On appelle cette condition la \emph{condition d'optimalité du second ordre}. 
\end{theorem}

\subsection{Conditions Suffisantes}

Passons maintenant aux conditions suffisantes d'optimalités... 

\begin{theorem}[CS1]
    Soit $f : \R^n \longrightarrow \R$ une fonction quelconque et $ \overline{x} \in \R^n$. On suppose qu'il existe 
    $r > 0 $ tel que : 
    \begin{itemize}
        \item $f$ est différentiable en tout point de $B(\overline{x}, r)$ 
        \item $ \forall x \in B( \overline{x}, r) / \{\overline{x}\}, \; \langle \nabla f(x), x - \overline{x} \rangle \geqslant 0 $
    \end{itemize}
    Alors $ \overline{x}$ est un \emph{minimum local de $f$}. 
\end{theorem}

\begin{theorem}[CS2]
    Soit $f : \R^n \longrightarrow \R$ une fonction quelconque et $ \overline{x} \in \R^n$. Alors si :
    \begin{itemize}
        \item $f$ est deux fois différentiable en $ \overline{x}$ 
        \item $ \nabla f( \overline{x}) = 0$ 
        \item $ \nabla^2 f( \overline{x}) \succeq 0 $
    \end{itemize}
    Alors $ \overline{x}$ est un \emph{minimum local strict de $f$}. 
\end{theorem}

Le théorème suivant nous permet de généraliser ce minimum local en minimum global. Pour cela, 
il faut que la fonction ne décroit plus (i.e que $ \overline{x}$ soit en bas de la parabole). 
Pour cela, on utilisera la notion de convexité. 

\begin{theorem}[Convexité et Minimum]
    Soit $f : \R^n \longrightarrow \R$ une fonction quelconque et $ \overline{x} \in \R^n $ un minimum 
    local de $f$. Alors : 
    \begin{enumerate}
        \item Si $f$ est \emph{convexe} alors $ \overline{x}$ est un minimum \emph{global} de $f$. 
        \item Si $f$ est \emph{strictement convexe} alors $ \overline{x}$ est \emph{l'unique minimum global} de $f$. 
    \end{enumerate}
\end{theorem}

% ========================================================================================================
% Algorithmes de descente

\section{Algorithmes de Descente}

\subsection{Principe}

Nous allons ici nous intéresser à des algorithmes permettant d'approximer un minimum d'un champ de scalaires. 
Ces algorithmes seront itératifs, autrement dit, à chaque pas de temps, nous approcherons de plus en plus 
notre solution. On les appelle des \emph{algorithmes de descente}. 
La notion de "descente" se formalise par la construction d'une suite $(u_n)_{n \in \N}$ tel que : 
    $$ \forall n \in \N, \quad f(u_{n+1}) \leqslant f(u_n) $$ 

\begin{definition}[Direction de descente]
    Soit $f : \R^n \longrightarrow \R$ une application différentiable et $x,d \in \R^n$. 
    Le vecteur $d$ est \emph{la direction de descente de $f$ en $x$} si : 
        $$ \langle \nabla f(x), d \rangle < 0 $$ 
\end{definition}

Pour l'algorithme de descente, ce vecteur $d$ nous donnera le "sens" dans lequel il faut descendre 
pour s'approcher du minimum. Or maintenant, que l'on sait de quel "coté" descendre, essayons de 
savoir "de combien" faut-il descendre. 

\begin{theorem}[Existence du pas]
    Soit $f : \R^n \longrightarrow \R$ une application différentiable. Soient $x,d \in \R^n$ tels que 
    $ \nabla f(x) \not = 0$. Si $d$ est une direction de descente de $f$ en $x$, alors il existe $ \mu > 0$ 
    tel que : 
        $$ \boxed{\forall \alpha \in ]0, \mu[, \quad f(x + \alpha d) < f(x)} $$ 
    De plus, pour tout $ \beta < 1$, il existe $ \mu > 0$ tel que : 
        $$ \boxed{\forall \alpha \in ]0, \mu[, \quad f(x + \alpha d) < f(x) \alpha \beta \langle \nabla f(x), d \rangle} $$ 
    On appelera $ \alpha$ le \emph{pas} de l'itération. 
\end{theorem}

\begin{proposition}
    Chaque itération du processus de descente se fait donc en 3 étapes : 
    \begin{enumerate}
        \item Trouver une direction $d_k$ telle que $ \langle \nabla f(x_k) , d_k \rangle < 0$ 
        \item Trouver un pas $ \alpha_k$ tel que $f(x_k + \alpha_k d_k) < f(x_k)$
        \item Calculer $x_{k+1} = x_k + \alpha_k d_k $
    \end{enumerate}
\end{proposition}

\subsection{Plus profonde descente}

Nous allons maintenant nous intéresser à différents types de descente. En effet, en fonction de l'application 
à minimiser nous utiliserons différents pas et directions lors de la descente. 
La méthode naturelle serait de prendre la direction dans laquelle les valeurs de l'application descendent le plus. 
Nous choissirons donc au début la direction $ d = - \nabla f(x)$. 

\begin{theorem}[Gradient comme direction]
    Soit $f : \R^n \longrightarrow \R$ une application différentiable. Soient $x \in \R^n$ alors 
    pour tout $d \in \R^n$ tel que $ ||d|| = || - \nabla f(x) ||$ on a : 
        $$ \langle \nabla f(x), - \nabla f(x) \rangle \leqslant \langle \nabla f(x), d \rangle $$ 
    On diras alors que $d^* = - \nabla f(x)$ est \emph{la direction de plus profonde descente}. 
\end{theorem}

\begin{proposition}
    Pour mettre en oeuvre notre plus profonde descente, à chaque itération, nous calculons donc : 
        $$ x_{k+1} = x_k - \alpha_k \nabla f(x_k) $$ 
    Un théorème précédent nous assure de pouvoir trouver un $ \alpha_k$ qui convient 
    à chaque itération. Pour une valeur de $ \alpha_k$ optimale, on pourrait résoudre le problème suivant à chaque itération : 
        \begin{mini*}|s|
            {}{f(x_k + \alpha d_k )}{}{}
            % \addConstraint{g_{j}(x)}{\le 0, j = 1, \hdots, m}
            \addConstraint{        \alpha }{> 0}
        \end{mini*}
    Mais nous ne sommes pas sûr de toujours pouvoir le résoudre. 
\end{proposition}

\subsection{Recherche Linéaire : Conditions}

On remarque assez vite que choisir un pas constant très petit (comme $0.0001$) rallonge considérablement 
la descente. D'autre part, nous savons que l'existence d'un minimum pour un champ de scalaire est une notion 
locale. Prendre un pas trop grand risquerait donc de nous faire rater le minimum et de faire augmenter la fonction 
objectif. 

Or, résoudre un problème d'optimisation pour $ \alpha_k$ à chaque itération 
n'est pas non plus une bonne solution car elle augmente la complexité de la méthode. 
Nous allons donc essayer de développer une méthode qui permet de résoudre approximativement la recherche 
d'un $ \alpha_k$ à moindre coût. Pour cela, nous choisirons $ \alpha_k$ en fonction de la rapidité ou non de 
la décroissance de notre fonction objectif. 

\begin{proposition}
    Pour trouver un pas raisonnable à chaque itération, nous utiliserons l'information contenue dans le gradient 
    de la fonction objectif. En effet, celui-ci nous permet de savoir comment celle-ci diminue et à quelle vitesse. 
    On regardera, pour chaque pas possible, la \emph{diminution} qu'il produit dans notre algorithme : 
        $$ f(x_k) - f(x_k + \alpha_k d_k) > \gamma \alpha_k, \quad \text{pour } \gamma > 0  $$ 
    Diminution qui sera proportionnelle au pas. 
    Ainsi, plus la pente $ \langle \nabla f(x_k), d \rangle $ de la fonction objectif dans la direction de $d$ 
    sera importante, plus le pas sera grand.   
\end{proposition}

\begin{definition}[Première condition de Wolfe - Condition d'Arjimo]
    Soit $f : \R^n \longrightarrow \R$ une application différentiable. Soient $x_k \in \R^n$ 
    et $d_k \in \R^n$ une direction de descente de $f$ en $x_k$ et $ \alpha_k > 0$. 
    On dit que la fonction $f$ \emph{diminue suffisament} en $x_k + \alpha_k d_k$ par rapport à $x_k$ si : 
        $$ {f(x_k + \alpha_k d_k) < f(x_k) + \alpha_k \beta \langle \nabla f(x_k), d_k \rangle} $$
    avec $ \beta \in ]0,1[$. 
\end{definition}

Cette condition nous permet de rejeter des pas trop grands qui ne fournissent pas de décroissance suffisante de 
la fonction. 

Il peut arriver que les pas $ \alpha_k$ deviennent très proches de 0 et empêchent la méthode de progresser. 
On va donc chercher à avancer dans notre descente mais sans dépasser le minimum (première condition de Wolfe), mais 
en avançant suffisament pour arriver sur un point où la pente de la fonction objectif est moins raide (seconde condition 
de Wolfe). 

\begin{definition}[Seconde Condition de Wolfe]
    Soit $f : \R^n \longrightarrow \R$ une application différentiable. Soit $x_k \in \R^n$, 
    $d_k \in \R^n$ une direction de descente de $f$ en $x_k$ et $ \alpha_k > 0$. 
    On dira que le point $x_k + \alpha_k d_k$ \emph{apporte une progression suffisante} par rapport 
    à $x_k$ si : 
        $$ \langle \nabla f(x_k + \alpha_k d_k), d_k \rangle \geqslant \beta \langle \nabla f(x_k), d_k \rangle $$ 
        $$ \iff \frac{ \langle \nabla f(x_k + \alpha_k d_k), d_k \rangle}{\langle \nabla f(x_k), d_k \rangle} \geqslant \beta $$ 
\end{definition}

Les algorithmes de \emph{recherche linéaire} nous permettrons d'identifier un pas qui satisfait les deux conditions de Wolfe. 

\subsection{Algorithme de Recherche Linéaire}

L'idée de l'algorithme de recherche linéaire est de construire de manière itérative un intervalle $[ \alpha_g, \alpha_d]$ 
contenant des valeurs qui satisfassent les deux conditions de Wolfe. On choisis ensuite un $ \alpha \in [ \alpha_g, \alpha_d]$
et s'il viole l'une des deux conditions de Wolfe, on réduit l'intervalle précédente et on recommence. 

% \begin{proposition}
%     Soient $f : \R^n \longrightarrow \R$ une application différentiable, $x_k \in \R^n$, une direction de 
%     descente $d_k \in \R^n$ de $f$ par rapport à $x_k$ et $ \alpha_0 > 0$. Soient 
%     $ \beta_1$ et $ \beta_2$ deux paramètres tels que $ 0 < \beta_1 < \beta_2 < 1$ et $ \lambda > 1$. 
%     On a l'algorithme suivant : 
%     \begin{enumerate}
%         \item \textbf{Initialisation} : $ i = 0, \alpha_g = 0, \alpha_d = \infty$ 
%         \item \textbf{Itération} : Tant que l'un des conditions de Wolfe est violée avec $ \alpha_k = \alpha_i$ faire : 
%             \begin{enumerate}
%                 \item si $ \alpha_i$ viole la première condition de Wolfe, alors le pas est trop long donc : 
%                     $$ \alpha_d = \alpha_i $$ 
%                     $$ \alpha_{i+1} = \frac{ \alpha_g + \alpha_d}{2} $$ 
%                 \item si $ \alpha_i$ viole la seconde condition de Wolfe, alors le pas est trop court donc : 
%                     $$ \alpha_g = \alpha_i $$ 
%                     $$ \alpha_{i+1} = 
%                         \begin{cases}
%                             \frac{ \alpha_g + \alpha_d}{2} \quad \text{si } \alpha_d < \infty \\ 
%                             \lambda \alpha_i \quad \text{sinon}  
%                         \end{cases}
%                     $$ 
%                 \item $ i++$ 
%             \end{enumerate}
%     \end{enumerate}
% \end{proposition}

\begin{algorithm}[H]
    \caption{Recherche linéaire satisfaisant les conditions de Wolfe}
        \begin{algorithmic}[1]
        \Require Fonction $f : \R^n \to \R$, point courant $x_k$, direction de descente $d_k$, pas initial $\alpha_0 > 0$, paramètres $0 < \beta_1 < \beta_2 < 1$, $\lambda > 1$
        \Ensure Pas $\alpha$ satisfaisant les conditions de Wolfe
        \State $i \gets 0$, $\alpha_g \gets 0$, $\alpha_d \gets \infty$
        \State $\alpha_i \gets \alpha_0$
        \While{$\alpha_i$ viole l'une des conditions de Wolfe}
            \If{$\alpha_i$ viole la première condition de Wolfe (Armijo)}
                \State $\alpha_d \gets \alpha_i$ \Comment{le pas est trop long}
                \State $\alpha_{i+1} \gets \frac{\alpha_g + \alpha_d}{2}$
            \ElsIf{$\alpha_i$ viole la seconde condition de Wolfe (curvature)}
                \State $\alpha_g \gets \alpha_i$ \Comment{le pas est trop court}
                \If{$\alpha_d < \infty$}
                    \State $\alpha_{i+1} \gets \frac{\alpha_g + \alpha_d}{2}$
                \Else
                    \State $\alpha_{i+1} \gets \lambda \alpha_i$
                \EndIf
            \EndIf
            \State $i \gets i + 1$
            \State $\alpha_i \gets \alpha_{i+1}$
        \EndWhile
        \State \Return $\alpha_i$
    \end{algorithmic}
\end{algorithm}


\begin{theorem}[Finitude de la recherche linéaire]
    Soit $f$ une application différentiable bornée inférieurement dans une direction $d_k$ 
    par rapport à $x_k \in \R^n$. Alors la recherche linéaire d'un pas se termine après un nombre 
    fini d'itérations. 
\end{theorem}

\subsection{Méthode de Newton}

Le principe de la méthode de Newton est d'approximer localement la fonction $f$ à minimiser 
par un développement de Taylor à l'ordre 2 puis de trouver le minimum de cette approximation. 

Dans sa forme la plus générale, la méthode de Newton sert à résoudre un système non linéaire 
définit par $g(x) = 0$ où $ g : \R^n \longrightarrow \R^n$. Elle consiste 
à construire une suite de points $(x_k)_{n \in \N}$ où à chaque étape on cherche un $d \in \R^n$ tel que :
    $$ \forall k \in \N, \quad g(x_k + d) = 0 $$ 
Or, en remplaçcant $g$ par son approximation linéaire : 
    $$ g(x_k + d) \simeq g(x_k) + J(x_k)d $$ 
La recherche de $d$ tel que $g(x_k + d) = 0$ est alors équivalente à la recherche d'un $d$ tel que : 
    \begin{align*}
        & g(x_k) + J(x_k)d = 0 \\ 
        \iff & g(x_k) = - J(x_k)d 
    \end{align*}
L'itéré suivant est alors $x_{k+1} = x_k + d$. On s'arrête lorsque $\| g(x_k) \|$ est plus petit qu'un seuil fixé. 

\begin{proposition}
    Appliquons maintenant la méthode de Newton à la recherche d'un minimum d'une fonction $f : \R^n \longrightarrow \R$. 
    On cherche donc à obtenir une condition d'optimalité de la forme $g(x) = \nabla f(x) = 0$. Dans le 
    cadre précédent, on a donc $g(x) = \nabla f(x)$ et $ J(x) = H_f(x)$. Ici, l'approximation de $f$ 
    par un développement de Taylor de la forme : 
        $$ f(x) \simeq f(x_k) + \langle \nabla f(x_k), d \rangle + \frac{1}{2} \langle d, H_f(x_k) d \rangle $$ 
    est minimisée si $ H_f(x_k)$ est définie positive. Autrement dit, lorsque le gradient de $g$ est nul soit : 
        $$ \nabla f(x_k) + H_f(x_k)d = 0 $$ 
\end{proposition}

On utilisera donc l'algorithme suivant : 

\begin{algorithm}[H]
    \caption{Méthode de Newton (optimisation)}
    \begin{algorithmic}[1]
        \Require Fonction $f : \mathbb{R}^n \to \mathbb{R}$, point initial $x_0$, tolérance $\varepsilon > 0$
        \Ensure Approximation $x^*$ d’un point critique de $f$
        \State $k \gets 0$
        \While{$\| \nabla f(x_k) \| > \varepsilon$}
            \State Calculer le gradient $\nabla f(x_k)$
            \State Calculer la hessienne $H_f(x_k)$
            \State Résoudre $H_f(x_k) d_k = - \nabla f(x_k)$
            \State Mettre à jour $x_{k+1} \gets x_k + d_k$
            \State $k \gets k + 1$
        \EndWhile
        \State \Return $x^* \gets x_k$
    \end{algorithmic}
\end{algorithm}


L'avantage de la méthode de Newton est sa convergence quadratique lorsque les conditions favorables sont réunies. 
En revanche, elle possède plusieurs inconvénients qui nous pousseront à développer une méthode plus fine : 
\begin{itemize}
    \item Si le point de départ de la méthode est trop éloigné, elle peut diverger très violement. 
    \item Elle nécessite le calcul de la hessienne de la fonction à minimiser. 
    \item Elle n'est pas définie si la hesienne n'est pas inversible. 
    \item Dans le contexte de l’optimisation, même lorsque la suite $(x_k)_{k \in \N}$
    converge, elle converge vers un point critique, c’est-à-dire une solution de l’équation 
    $ \nabla f(x) = 0$. Elle peut donc converger vers un maximum local.
\end{itemize}

Pour pallier ces contraintes, nous utiliserons la méthode de Newton hybride. 
Elle vise à améliorer la robustesse de la méthode de Newton classique, 
en garantissant que la direction de recherche reste toujours une direction de descente. 
Elle combine pour cela la régularisation de la hessienne et une étape de recherche linéaire. 
Ainsi, même lorsque la hessienne $H_f(x_k)$ n’est pas définie positive, la méthode reste stable et converge vers un minimum local.

\begin{algorithm}[H]
    \caption{Méthode de Newton hybride (régularisée avec recherche linéaire)}
    \begin{algorithmic}[1]
        \Require Fonction $f : \mathbb{R}^n \to \mathbb{R}$, point initial $x_0$, tolérance $\varepsilon > 0$
        \Ensure Approximation $x^*$ d’un point critique de $f$
        \State $k \gets 0$
        \While{$\| \nabla f(x_k) \| > \varepsilon$}
            \State Calculer le gradient $\nabla f(x_k)$ et la hessienne $H_f(x_k)$
            \If{$H_f(x_k)$ est définie positive}
                \State $D_k \gets H_f(x_k)$
            \Else
                \State Choisir $\mu > 0$ tel que $D_k = H_f(x_k) + \mu I$ soit définie positive
            \EndIf
            \State Résoudre $D_k d_k = - \nabla f(x_k)$
            \State Effectuer une recherche linéaire pour trouver $\alpha_k > 0$ satisfaisant la condition de Wolfe ou d’Armijo
            \State Mettre à jour $x_{k+1} \gets x_k + \alpha_k d_k$
            \State $k \gets k + 1$
        \EndWhile
        \State \Return $x^* \gets x_k$
    \end{algorithmic}
\end{algorithm}

\subsection{Méthode de quasi-Newton (BFGS)}

Comme dit précédement, la méthode de Newton possède quelques désavantages que l'on souhaiterai résoudre. 
En effet, la hessienne est parfois impossible à calculer si la fonction n'est pas deux fois dérivable. De plus, 
une fois calculée, il peut être très coûteux de l'évaluer en un point. Nous allons donc ici chercher 
à estimer trouver une estimation $M$ de cette hessienne ou de son inverse (notée $W$).

Le schéma d'un algorithme de quasi-Newton est le même que la méthode de Newton hybride dans lequel on 
remplace la résolution du système linéaire : 
    $$ H_f(x_k) d_k = - \nabla f(x_k) $$ 
par la résolution d'un autre système : 
    $$ M_k d_k = - \nabla f(x_k) \quad \text{ou} \quad d_k = - W_k \nabla f(x_k)$$ 
où $M_k$ ou $W_k$ sont des matrices construires grâce à des \emph{formules de mise à jour}. 

\begin{proposition}
    Soient 
        \[ s_k = x_{k+1} - x_k, \quad y_k = \nabla f(x_{k+1}) - \nabla f(x_k) \] 
    la mise à jour de $W_k$ pour obtenir $W_{k+1}$ est donnée par :

    \[ \boxed{ W_{k+1} = \Big(I - \frac{s_k y_k^\top}{y_k^\top s_k}\Big) W_k \Big(I - \frac{y_k s_k^\top}{y_k^\top s_k}\Big) + \frac{s_k s_k^\top}{y_k^\top s_k}} \]

    Cette formule garantit que $W_{k+1}$ reste symétrique et définie positive à condition que $y_k^\top s_k > 0$.
    La direction de descente est alors simplement :
        \[ \boxed{d_k = - W_k \nabla f(x_k)} \]
\end{proposition}

\begin{algorithm}[H]
\caption{Algorithme quasi-Newton (BFGS)}
\begin{algorithmic}[1]
\Require Fonction $f : \mathbb{R}^n \to \mathbb{R}$, point initial $x_0$, tolérance $\varepsilon > 0$, matrice initiale $W_0$ définie positive
\Ensure Approximation $x^*$ d’un minimum local de $f$
\State $k \gets 0$
\While{$\| \nabla f(x_k) \| > \varepsilon$}
    \State Calculer le gradient $g_k \gets \nabla f(x_k)$
    \State Déterminer la direction de descente $d_k \gets - W_k g_k$
    \State Effectuer une recherche linéaire pour obtenir $\alpha_k > 0$
    \State Mettre à jour $x_{k+1} \gets x_k + \alpha_k d_k$
    \State Calculer $s_k \gets x_{k+1} - x_k$, $y_k \gets \nabla f(x_{k+1}) - \nabla f(x_k)$
    \State Mettre à jour la matrice inverse de Hessienne :
    \[
        W_{k+1} \gets \Big(I - \frac{s_k y_k^\top}{y_k^\top s_k}\Big) W_k \Big(I - \frac{y_k s_k^\top}{y_k^\top s_k}\Big) + \frac{s_k s_k^\top}{y_k^\top s_k}
    \]
    \State $k \gets k + 1$
\EndWhile
\State \Return $x^* \gets x_k$
\end{algorithmic}
\end{algorithm}

